\section{Data and Benchmark}
\label{sec:data}

\subsection{Legal Data Sources}

The corpus was constructed from the 2015 Vietnamese Penal Code and the
relevant amendments and supplements
\cite{quochoi2015blhs,quochoi2017suadoiblhs}. The data source focuses on
provisions that describe conditions of applicability and legal effects in
criminal law. The study does not use judgments, case files, personal data,
or empirical behavioral data.

The input unit of the pipeline is not an entire legal article represented as
a single passage, but rather a set of segmented and normalized rules. A legal
article may yield multiple rules if it contains several distinct
condition--effect pairs. Conversely, multiple legal articles may lead to the
same effect event while applying to different conditions or legal subjects.

Each rule record contains at least the following fields:

\begin{itemize}
    \item \texttt{rule\_id}: the unique identifier of the rule;
    \item \texttt{article\_id}: the legal article number;
    \item \texttt{article\_title}: the title of the legal article;
    \item \texttt{legal\_subject}: the legal subject;
    \item \texttt{condition}: the condition of applicability;
    \item \texttt{effect}: the legal effect;
    \item \texttt{condition\_event}: the condition event identifier;
    \item \texttt{condition\_event\_name}: the human-readable name of the
    condition event;
    \item \texttt{effect\_event}: the effect event identifier;
    \item \texttt{effect\_event\_name}: the human-readable name of the effect
    event;
    \item \texttt{causal\_type}: the relation type or effect type, when
    available.
\end{itemize}

The legal corpus is used as a source of normative knowledge. The system does
not infer that the condition--effect relations in the corpus represent
statistical causal relations in the empirical world. All conclusions produced
by the pipeline are bounded by the content of the encoded rules.

\subsection{Rule Normalization Procedure}

The normalization process converts legal text into a unified schema for graph
construction, vector indexing, and evaluation. The procedure consists of the
following main steps:

\begin{enumerate}
    \item identify the legal article and its corresponding title;

    \item segment the content into a condition of applicability and a legal
    effect;

    \item identify the subject affected by the rule;

    \item normalize the condition and effect into concise event names;

    \item create a stable event ID from the event name;

    \item assign a rule ID and link the rule to an article ID;

    \item validate the required fields and remove records with insufficient
    information.
\end{enumerate}

An event ID is normalized by converting the content to lowercase, removing
Vietnamese diacritics, replacing spaces with underscores, and removing
unnecessary characters. For example:

\begin{quote}
Event name: ``Phạm tội chưa đạt'';

Event ID: \texttt{pham\_toi\_chua\_dat}.
\end{quote}

Different legal expressions with the same meaning may be mapped to the same
event ID. Conversely, effects that differ materially in penalty level,
applicable subject, or condition are retained as distinct events. This
normalization reduces surface-form duplication while preserving the
distinctions necessary for legal reasoning.

The current procedure does not fully normalize every exception, negation,
cross-reference, or multi-condition structure. Most rules in the experimental
corpus are represented by one positive condition event and one positive effect
event. The corpus is therefore most suitable for reasoning over short
condition--effect chains.

\subsection{Corpus and Graph Statistics}

The normalized corpus contains \(2\,884\) rules and \(1\,814\) events. The
heterogeneous graph constructed from the corpus contains \(5\,134\) nodes and
\(14\,420\) edges. The main statistics are presented in
Table~\ref{tab:graph}.

\begin{table}[t]
    \centering
    \caption{Statistics of the legal corpus and causal graph.}
    \label{tab:graph}
    \begin{tabular}{lr}
        \hline
        \textbf{Component} & \textbf{Count} \\
        \hline
        Legal rules & \(2\,884\) \\
        Legal events & \(1\,814\) \\
        Total nodes & \(5\,134\) \\
        Total edges & \(14\,420\) \\
        \texttt{CAUSES} edges & \(2\,884\) \\
        Unique directed event pairs & \(2\,122\) \\
        Two-hop causal chains & \(112\) \\
        \hline
    \end{tabular}
\end{table}

The number of \texttt{CAUSES} edges equals the number of rules because each
rule creates at least one relation from a condition event to an effect event.
However, the number of unique event pairs is lower than the number of
\texttt{CAUSES} edges, indicating that some event pairs are supported by
multiple rules or multiple legal grounds.

The total number of edges is substantially larger than the number of
\texttt{CAUSES} edges because the graph also stores heterogeneous links among
events, rules, and articles. These edges support evidence provenance but are
not treated as causal edges during path traversal.

From the event--event projection, the system identifies \(112\) two-hop
causal chains of the form:

\[
v_0 \rightarrow v_1 \rightarrow v_2.
\]

A chain may be used to generate multiple questions with different answer
intents, such as asking for the outcome, the cause, the mediator, or an
explanation of the complete path. Consequently, the number of benchmark
questions may exceed the number of unique two-hop chains.

\subsection{Benchmark Construction Objectives}

The benchmark was constructed to separately evaluate the following
capabilities:

\begin{enumerate}
    \item retrieving the correct rules and articles;
    \item retrieving the correct events;
    \item recovering the causal path in the correct order;
    \item distinguishing a direct relation from a relation mediated by an
    intermediate event;
    \item verifying the claim expressed in the question;
    \item generating an answer that satisfies the answer intent;
    \item citing the correct legal grounds.
\end{enumerate}

The benchmark design emphasizes evaluation of the supporting path rather than
only the final answer. This approach is consistent with the principle of
evaluating supporting facts in multi-hop question answering
\cite{yang2018hotpotqa}, but is adapted to the rule and event structure of the
legal corpus.

The benchmark is not intended to replace general legal evaluation suites such
as LegalBench \cite{guha2023legalbench}. It is used as a controlled testbed for
evaluating the pipeline on the specific type of legal causal structure that
the system is designed to process.

\subsection{Question Generation Procedure}

Questions are generated from causal paths, rules, and articles contained in
the graph. Given a two-hop path:

\[
A \xrightarrow{r_1} B \xrightarrow{r_2} C,
\]

the system can generate the following question types:

\begin{description}
    \item[Forward multihop] Given \(A\), identify the final effect \(C\).

    \item[Reverse reasoning] Given \(C\) or a sequence of results, identify
    the initial cause \(A\).

    \item[Bridge event] Given \(A\) and \(C\), identify the mediator \(B\).

    \item[Causal chain explanation] Present the complete
    \(A \rightarrow B \rightarrow C\) chain and the corresponding legal
    grounds.

    \item[Yes/no positive] Determine whether the stated relation or chain is
    supported by the graph.

    \item[Yes/no counterexample] Assert a direct relation
    \(A \rightarrow C\) when the graph supports only the indirect path
    \(A \rightarrow B \rightarrow C\).

    \item[Branch reasoning] Aggregate multiple outcomes that share a common
    prefix or mediator.

    \item[Convergence reasoning] Aggregate multiple initial conditions that
    lead to a common suffix or outcome.

    \item[Bridge-convergence fallback] Use a convergent structure when the
    requirements for generating a complete convergence sample are not met,
    while still allowing a question about the bridge event or main path to be
    generated.
\end{description}

The questions are generated from controlled templates and use human-readable
event names rather than event IDs. Templates make it possible to define the
gold answer and gold path consistently, but they may also produce repetitive
linguistic patterns. This is a threat to generalizability and is treated as a
benchmark limitation.

\subsection{Question-Type Distribution}

The benchmark contains a total of \(250\) questions. The observed distribution
is presented in Table~\ref{tab:benchmark}.

\begin{table}[t]
    \centering
    \caption{Distribution of question types in the benchmark.}
    \label{tab:benchmark}
    \begin{tabular}{lr}
        \hline
        \textbf{Question type} & \textbf{Count} \\
        \hline
        Forward multihop & \(70\) \\
        Reverse reasoning & \(45\) \\
        Bridge event & \(40\) \\
        Causal chain explanation & \(35\) \\
        Yes/no positive & \(20\) \\
        Yes/no counterexample & \(20\) \\
        Branch reasoning & \(4\) \\
        Convergence reasoning & \(10\) \\
        Bridge-convergence fallback & \(6\) \\
        \hline
        \textbf{Total} & \(\mathbf{250}\) \\
        \hline
    \end{tabular}
\end{table}

The four categories of forward multihop, reverse reasoning, bridge event, and
causal chain explanation constitute most of the benchmark. This distribution
is consistent with the primary objective of evaluating reasoning over two-hop
chains. Branch and convergence questions are less frequent because the graph
contains only a limited number of structures satisfying the requirements for
multiple paths with a common prefix or suffix.

\subsection{Ground Truth and Evaluation Schema}

Each benchmark sample has a unique identifier and contains two groups of
fields: fields describing the question and fields used for evaluation.

The main descriptive fields include:

\begin{itemize}
    \item \texttt{id};
    \item \texttt{question};
    \item \texttt{question\_type};
    \item metadata about the source path or structure used to generate the
    question.
\end{itemize}

The \texttt{evaluation} object contains ground truth for the different
layers of evaluation:

\begin{description}
    \item[Retrieval ground truth]
    Includes \texttt{gold\_rule\_ids},
    \texttt{gold\_article\_ids}, and the expected retrieval size.

    \item[Event ground truth]
    Includes \texttt{gold\_event\_ids}, event names, and the final event.

    \item[Path ground truth]
    Includes the ordered steps of the gold path as source event--target event
    pairs.

    \item[Verification ground truth]
    Includes \texttt{requires\_counterfactual},
    \texttt{gold\_decision}, and the verification target.

    \item[Answer ground truth]
    Contains the reference answer appropriate to the question type.

    \item[Citation ground truth]
    Contains the article IDs expected to support the answer.

    \item[Difficulty metadata]
    Contains retrieval, reasoning, and generation difficulty, together with an
    overall score.
\end{description}

All \(250\) questions have an evaluation object, gold rules, gold events, a
gold answer, and gold citations. A total of \(236\) questions have a clearly
defined linear gold path. The remaining branch and convergence questions
require evaluation over multiple paths or a subgraph rather than a single
linear chain.

\subsection{Gold Verification Decisions}

The distribution of gold decisions is as follows:

\begin{table}[t]
    \centering
    \caption{Distribution of gold verification decisions.}
    \label{tab:gold-decision}
    \begin{tabular}{lr}
        \hline
        \textbf{Gold decision} & \textbf{Count} \\
        \hline
        \texttt{SUPPORTED} & \(230\) \\
        \texttt{REJECT\_DIRECT\_CLAIM} & \(20\) \\
        \texttt{UNCERTAIN} & \(0\) \\
        \hline
        \textbf{Total} & \(\mathbf{250}\) \\
        \hline
    \end{tabular}
\end{table}

The benchmark is strongly imbalanced with respect to decision labels: \(92\%\)
of the samples belong to the \texttt{SUPPORTED} class, whereas only \(8\%\)
belong to the \texttt{REJECT\_DIRECT\_CLAIM} class. Verification accuracy is
therefore insufficient as a standalone evaluation measure. The study also
reports Macro-F1, Balanced Accuracy, and a confusion matrix.

The benchmark contains no gold sample in the \texttt{UNCERTAIN} class.
Consequently, the ability to detect insufficient evidence or an indeterminate
state is not directly evaluated through a separate gold class.

\subsection{Counterfactual Subset}

The benchmark marks \(20\) questions as requiring counterfactual verification.
However, content inspection shows that these samples are primarily yes/no
counterexamples concerning direct relations. Given the path:

\[
A \rightarrow B \rightarrow C,
\]

a question may assert that \(A\) directly leads to \(C\). The gold decision is
then \texttt{REJECT\_DIRECT\_CLAIM} because the graph contains no direct edge
\(A \rightarrow C\).

These samples effectively evaluate the ability to distinguish a direct edge
from an indirect path, but they do not provide complete ground truth for the
intervention:

\[
do(B=\mathrm{FALSE}).
\]

Specifically, the benchmark does not yet include expert annotations for:

\begin{itemize}
    \item the factual state of the mediator;
    \item the factual state of the outcome;
    \item the counterfactual outcome after intervention;
    \item the rules disabled by the intervention;
    \item alternative paths after intervention;
    \item whether the mediator is genuinely necessary in the legal model.
\end{itemize}

Accordingly, the ``counterfactual verification accuracy'' metric in the
current study primarily reflects the ability to process negative direct-edge
claims. It is not interpreted as evidence that LegalSCM has been validated on
a comprehensive hard-intervention benchmark.

\subsection{Benchmark Difficulty}

The benchmark assigns difficulty metadata based on the retrieval, reasoning,
and generation characteristics of each sample. In the current dataset:

\begin{itemize}
    \item \(181\) questions are labeled \texttt{medium};
    \item \(69\) questions are labeled \texttt{hard};
    \item no question is labeled \texttt{easy}.
\end{itemize}

Hard samples typically require one of the following characteristics:

\begin{itemize}
    \item multiple supporting rules or articles;
    \item multiple candidate paths sharing a mediator or outcome;
    \item a branch or convergence structure;
    \item distinguishing a direct edge from an indirect path;
    \item an answer that aggregates multiple components.
\end{itemize}

The difficulty label is used for subgroup analysis and is not provided to the
retriever, verifier, or answer generator.

\subsection{Evaluation-Label Leakage Control}

Because the benchmark contains \texttt{question\_type}, gold paths, and gold
answers, these fields must be prevented from influencing prediction. The
pipeline applies the following principles:

\begin{enumerate}
    \item Step 3 receives only the \texttt{question} string and does not
    receive the evaluation object.

    \item Step 4 receives only the retrieval result and the query text stored
    in the artifact.

    \item The claim type is inferred from the query text and does not read
    \texttt{question\_type}.

    \item Step 5 infers the answer intent from the query text and the
    \texttt{query\_analysis} produced by Step 4.

    \item Gold rules, gold events, gold paths, gold decisions, and gold answers
    are used only by Step 7 after all predictions have been generated.

    \item The batch runner may copy \texttt{question\_type} into a prediction
    for subgroup analysis, but this field is not passed to the retrieval,
    verification, or answer-generation functions.
\end{enumerate}

This mechanism prevents the system from directly using the question type or
gold label to select an answer template. This is necessary for answer-intent
results to reflect query-analysis capability rather than leakage from the
benchmark.

\subsection{Integrity of the Paired Experiment}

The two principal runs use the same:

\begin{itemize}
    \item benchmark of \(250\) questions;
    \item graph, causal memory, embeddings, and FAISS index;
    \item retriever and all retrieval hyperparameters;
    \item query analyzer;
    \item Step 5 and the extractive answer provider;
    \item evaluator and metric set.
\end{itemize}

The only difference is the verification mode:

\begin{itemize}
    \item \texttt{structural\_scm};
    \item \texttt{path\_ablation}.
\end{itemize}

Both runs generate all \(250/250\) predictions and record no pipeline errors.
Holding all other components fixed makes this a valid paired ablation at the
implementation level.

Older artifacts generated by earlier versions of the runner or Step 5 are not
included in the paired comparison. Specifically, only files with the suffix
\texttt{\_v41\_step5\_v51} are used for the main results.

\subsection{Limitations and Threats to Benchmark Validity}

The current benchmark has several important limitations.

\paragraph{Shared graph source.}

The questions and ground truth are generated from the same graph used by the
pipeline. This provides control over the reasoning structure but may
overestimate the ability to recover encoded relations. The benchmark does not
yet thoroughly evaluate generalization to an independent corpus or graph.

\paragraph{Template effects.}

Many questions are generated from templates. Phrases such as ``final effect,''
``intermediate event,'' ``directly,'' or ``initial condition'' may help the
query analyzer identify the answer intent. Performance may be lower on natural
questions with more diverse formulations.

\paragraph{Lack of expert validation.}

The benchmark has not been comprehensively reviewed by legal experts at the
level of every question, path, and answer. The gold data reflect the structure
of the graph and normalized rules rather than official legal conclusions.

\paragraph{Limited counterfactual coverage.}

The counterfactual subset contains only \(20\) samples and primarily evaluates
negative direct-edge claims. The benchmark is not sufficiently sensitive to
distinguish hard intervention from node-deletion reachability.

\paragraph{Linear gold paths.}

The exact-path metric is based on a single linear chain, whereas branch and
convergence questions may require a set of paths or a subgraph. Consequently,
an exact-path score of zero for these groups does not necessarily imply an
incorrect answer.

\paragraph{Simplified rule representation.}

The corpus does not fully represent negation, exceptions, multiple concurrent
conditions, rule priority, or legal conflicts. The benchmark therefore does
not evaluate these capabilities either.

\paragraph{Narrow domain scope.}

The benchmark focuses exclusively on the Penal Code and is not representative
of other areas of law, such as civil, administrative, labor, or procedural
law.

In light of these limitations, the experimental results should be understood
as an internal evaluation on a controlled testbed. They must not be used to
conclude that the system can provide reliable legal advice in real-world
settings.
